\subsection{Two-way coupling}
\label{sec:coupling}

The two models exchange state through the MODFLOW~6 BMI/XMI application programming interface, which
exposes the solver's memory in place and runs the GWF6 flow and GWT6 transport models as a coupled
exchange. Surface-water percolation drives groundwater recharge; the groundwater--stream exchange
returns to the surface through the MODFLOW~6 streamflow-routing (SFR) package, whose reaches then
carry the discharged PFAS in-stream via the streamflow-transport (SFT) package (Section~\ref{sec:pfas}).

\paragraph{Surface to subsurface: recharge.} Because SWAT\textsuperscript{+} HRUs are
geo-located polygons, each HRU is intersected with the MODFLOW grid to build an area-weighted
HRU-to-cell map (no intermediate disaggregated-HRU step is required, unlike the embedded
SWAT--MODFLOW code). Per-cell recharge is the area-weighted sum of the percolation of the HRUs
overlapping that cell, divided by the cell area. The map is a sparse matrix, so a daily recharge
field is a single matrix--vector product. In coupled mode, the converged steady model is
rewritten as a sequence of daily stress periods (a steady spin-up period followed by transient
periods with specific storage and specific yield), and recharge is set at the start of each step
by writing in place to the MODFLOW recharge memory pointer through the BMI/XMI interface. On the
demonstration watershed all 1{,}096 daily steps of a three-year run converged, with simulated heads
tracking seasonal recharge within $\pm0.16$~m about the steady baseline.

\paragraph{Subsurface to surface: the groundwater--stream exchange.} The groundwater--stream
exchange is a native per-reach SFR budget term: the head-dependent leakage between each reach and the
cell it occupies, positive when groundwater discharges to the stream (gaining/baseflow) and negative
when the stream loses to the aquifer. SFR computes stage from routed flow and so enables in-stream
transport, whereas the simpler river (RIV) package prescribes a fixed stage per cell and can route
neither water nor solute. Because each reach is already a segment of a SWAT\textsuperscript{+}
channel, the exchange aggregates directly to the channel network with no spatial join.

The SFR reach network is built from the SWAT\textsuperscript{+} channel graph. Each channel
(filtered to stream order $\geq 2$) is oriented upstream-to-downstream by the DEM elevation of its
endpoints and intersected with the grid to give one reach per crossed active cell; reaches inherit
the channel width, a slope computed from the channel's elevation drop over its length, a streambed
top set below the land surface, and a streambed hydraulic conductivity. Connectivity follows a
single-successor graph (each reach has one downstream reach: the next cell within a channel, or the
first reach of the receiving channel at a confluence), with reaches numbered in topological order so
that every reach's downstream neighbour has a higher index --- matching the solver's preferred
ordering and avoiding a sign ambiguity in the zero-based connection list. On the demonstration
watershed the network has $1{,}506$ reaches across $273$ channels with no connectivity violations.
